\chapter{Swift算法设计与实现}\label{Chap:chap3}

本章详细介绍Swift算法的设计与实现，包括系统架构概述、决策问题建模与设计原则、11维观测空间设计、多信号融合拥塞检测算法、差异化拥塞响应机制、启发式反馈自适应的参数调优、融合控制策略以及每流独立状态管理机制。

\section{系统架构概述}

Swift算法的系统架构如图\ref{fig:architecture}所示，由四个核心模块组成：网络仿真模块、拥塞控制环境模块、智能决策模块和进程间通信模块。

\begin{figure}[htbp]
\centering
\begin{tikzpicture}[
    module/.style={rectangle, draw, thick, minimum width=4.5cm, minimum height=1.2cm, align=center, rounded corners=3pt, font=\small},
    submod/.style={rectangle, draw, thin, minimum width=3.5cm, minimum height=0.5cm, align=center, font=\scriptsize, fill=white},
    comm/.style={rectangle, draw, thick, fill=yellow!15, minimum width=1.8cm, minimum height=4.5cm, align=center, rounded corners=3pt, font=\small},
    arrow/.style={->, >=stealth, thick},
    darrow/.style={<->, >=stealth, thick},
]

% 网络仿真模块
\node[module, fill=blue!12] (sim) at (0,0) {\textbf{网络仿真模块}\\(ns-3 仿真器)};
\node[submod, below=0.1cm of sim] (sim1) {TCP/IP协议栈};
\node[submod, below=0.05cm of sim1] (sim2) {Dumbbell拓扑构建};
\node[submod, below=0.05cm of sim2] (sim3) {队列管理 (ECN/RED)};
\node[submod, below=0.05cm of sim3] (sim4) {FlowMonitor监控};

% 通信模块
\node[comm] (zmq) at (5.5, -0.8) {\textbf{IPC}\\\rotatebox{90}{ZeroMQ}\\~\\Req/Rep\\ProtoBuf};

% 智能决策模块
\node[module, fill=green!12] (agent) at (11,0) {\textbf{智能决策模块}\\(Python TcpSwift)};
\node[submod, below=0.1cm of agent] (a1) {多信号融合拥塞检测};
\node[submod, below=0.05cm of a1] (a2) {差异化拥塞响应};
\node[submod, below=0.05cm of a2] (a3) {自适应$\alpha$调优};
\node[submod, below=0.05cm of a3] (a4) {融合控制策略};

% 环境模块（底部跨越）
\node[module, fill=orange!12, minimum width=13cm] (env) at (5.5, -4.5) {\textbf{状态采集与决策应用模块}——11维状态采集 + 决策应用 + 反馈值计算};

% 箭头
\draw[arrow, blue!70] (sim3.east) -- ++(1.0,0) |- node[above, font=\scriptsize, pos=0.25] {11维观测} (zmq.west);
\draw[arrow, green!50!black] (zmq.east) -- ++(1.0,0) |- node[above, font=\scriptsize, pos=0.75] {$[ssThresh, cwnd]$} (a1.west);
\draw[arrow, red!70] (agent.south) -- ++(0,-0.3) -| node[right, font=\scriptsize, pos=0.3] {动作} (zmq.south east);
\draw[arrow, orange!70] (zmq.south west) -- ++(-1.0,0) |- node[left, font=\scriptsize, pos=0.3] {窗口更新} (sim2.east);

% 环境连接
\draw[darrow, gray] (sim4.south) -- ++(0,-0.5) -| (env.north west);
\draw[darrow, gray] (a4.south) -- ++(0,-0.5) -| (env.north east);

% 每流状态标注
\node[rectangle, draw, dashed, fill=gray!5, font=\scriptsize, minimum width=3cm, align=center] at (11, -3.2) {每流独立状态 (HashMap)\\socketUuid → FlowState};

\end{tikzpicture}
\caption{Swift系统架构}
\label{fig:architecture}
\end{figure}

\textbf{网络仿真模块}基于ns-3（Network Simulator 3）离散事件仿真平台构建，实现完整的TCP/IP协议栈。ns-3是目前学术界和工业界广泛使用的网络仿真工具，提供了丰富的网络协议模型和灵活的可扩展架构。网络仿真模块负责构建哑铃型等可控拓扑、模拟数据包传输过程、执行离散事件驱动的仿真调度，并通过FlowMonitor等工具记录吞吐量、时延、丢包和公平性指标。

\textbf{主动队列管理与ECN信号通路}。仿真器\path{contrib/opengym/examples/swift-tcp/sim.cc}将瓶颈两侧网关的默认队列规则设为\texttt{ns3::RedQueueDisc}，并按队列容量的比例配置标记门限：$MinTh = 0.3 \times Q$、$MaxTh = 0.9 \times Q$，其中$Q = \max(BDP \times n_{leaf} / MTU, 100)$为按带宽延迟积自动定尺的队列长度（单位为数据包）。同时置\texttt{RedQueueDisc::UseEcn}为真，使RED在队列长度进入$[MinTh, MaxTh]$区间时对ECT数据包置CE码点而非直接丢弃。

若瓶颈使用不设置CE码点的\texttt{PfifoFastQueueDisc}，则协议栈不能从该队列获得ECN标记，第~\ref{sec:multi-signal}节所述ECN分支也无法由瓶颈标记触发。因此本文的参考配置选用RED并开启ECN标记；第~\ref{Chap:chap4}章所引用的实验产物均在该配置下生成。\texttt{ns3::TcpSocketBase::UseEcn}被统一置为\texttt{On}并作用于四种被比较协议，以控制ECN协商能力这一实验变量。

Swift决策进程与ns-3仿真进程经ZeroMQ同步请求—应答路径交互：仿真进程传输包含4项元数据和11项有效状态的15元素容器，决策进程返回慢启动阈值与拥塞窗口决策；状态与动作的跨进程传输沿用ns3-gym/OpenGym框架提供的容器结构，仅作为消息通道，不含学习逻辑。由于该路径依赖同步消息交互，本文未将其与MTP多线程并行仿真机制结合使用。

\textbf{状态采集与决策应用模块}（环境层）负责在TCP协议栈与决策进程之间建立桥梁：五个TCP拥塞控制回调更新观测状态，其中\texttt{GetSsThresh()}与\texttt{IncreaseWindow()}触发同步决策交互；模块随后将返回的\texttt{ssthresh}与\texttt{cwnd}决策应用于TCP连接。

拥塞控制环境模块实现了以下回调接口：

\begin{itemize}
    \item \texttt{GetSsThresh()}：慢启动阈值获取回调，在协议栈请求窗口缩减参数时触发同步决策，具体原因需结合CA与ECN状态判定
    \item \texttt{IncreaseWindow()}：窗口增加回调，在收到正常确认时触发同步决策
    \item \texttt{PktsAcked()}：数据包确认回调，用于更新最近RTT与已确认报文段数量
    \item \texttt{CongestionStateSet()}：拥塞状态设置回调，用于处理状态迁移及超时时的陈旧动作作废
    \item \texttt{CwndEvent()}：拥塞窗口事件回调，用于更新ECN事件状态
\end{itemize}

\textbf{智能决策模块}采用Python实现，负责执行Swift算法的核心逻辑：多信号融合拥塞检测、差异化拥塞响应、自适应参数调优和融合控制策略计算。智能决策模块维护每个TCP连接的独立状态，包括历史度量缓冲区、拥塞事件计数、性能反馈状态和自适应参数。

\textbf{进程间通信模块}采用ZeroMQ（ZMQ）消息队列实现仿真进程和决策进程之间的同步通信。通信协议采用请求-应答（Request-Reply）模式：仿真进程在\texttt{GetSsThresh()}或\texttt{IncreaseWindow()}回调中发送观测，决策进程返回控制动作；其余回调用于更新下一次观测所需的环境状态。消息格式采用Protocol Buffers序列化。

\section{决策问题建模与设计原则}

Swift把网络拥塞控制刻画为\emph{逐事件决策问题}：在每个ACK确认或拥塞事件触发的回调上，决策模块以当前状态向量$s$为输入，按照确定性规则直接输出新的慢启动阈值与拥塞窗口决策$\langle ssthresh, cwnd\rangle$，并把由环境按传输进度与排队/丢包代价合成的性能反馈值$f_t$用于参数层的慢速自适应。该"状态—决策—反馈"三元结构不涉及策略学习：规则是离线设计的，反馈仅用于乘性增加因子等参数的小步调整，不存在训练、探索或模型更新环节。决策结构的合理性由以下四组设计原则保证。

\subsection{状态向量设计原则}

状态向量的设计决定了一次回调能够携带多少决策所需信息。良好设计的判定标准包括：

\textbf{状态信息覆盖}：状态应包含区分不同网络状况所需的信息。信息不足会使决策模块难以区分需要不同响应的情形，例如"ECN标记触发的收缩"与"普通丢包"。Swift的11维有效状态纳入ECN状态、拥塞控制状态机状态、拥塞事件、RTT与窗口度量，为三类拥塞语义判定提供输入，但不声称覆盖协议栈的全部可观测状态。

\textbf{可观测性（Observability）}：状态应能够从环境中直接获得。Swift的15个容器字段均由ns-3 TCP拥塞控制回调提供或由仿真环境生成，无需额外探测流量。

\textbf{维度可控性（Dimensional Controllability）}：状态维度应保持在规则可维护的范围内。Swift的11维有效状态配合基于规则的决策，单次处理涉及固定规模的字段读取、滑动窗口更新与比较；当前窗口容量固定，因此相对于仿真时长，其时间开销为常数级。

\subsection{决策输出设计}

决策输出决定了窗口控制的操作语义。Swift直接输出新的ssthresh与cwnd\emph{绝对值}而非窗口增量或离散档位，理由如下：

\textbf{控制粒度（Control Granularity）}：直接输出窗口绝对值可将目标窗口（如$\alpha \times BDP$）映射到协议状态，避免离散动作（如"增加10\%"）引入的额外量化。

\textbf{绝对值语义（Absolute Semantics）}：每步决策都以协议栈当前状态为基准重新计算；若待应用窗口因协议栈进入CA\_LOSS而失效，后续事件将基于更新后的状态重新计算。

\textbf{边界约束（Bound Constraints）}：Python控制器将拥塞窗口限制在$[4 \times MSS,\ \max(4 \times BDP,\ 200 \times MSS)]$区间；环境侧在应用动作前另以$2\times MSS$校验拥塞窗口和慢启动阈值的下界。二者共同限制当前实现中的窗口取值，但不替代对跨进程输入合法性的完整验证。

\subsection{性能反馈值设计}

环境侧在每个事件上合成一个标量性能反馈值$f_t$，供参数层的启发式自适应使用。其设计遵循三条原则：

\textbf{即时性（Immediacy）}：每个确认事件与每个窗口收缩回调都产生反馈，不存在稀疏反馈造成的信号缺失问题。

\textbf{多目标折中（Multi-objective Trade-off）}：反馈由传输进度激励与代价惩罚两部分合成：正常确认事件上的反馈为"与确认段数成正比且有上限的吞吐激励"减去"RTT膨胀惩罚"；窗口收缩回调上按拥塞类型施加分级负向惩罚（ECN、丢包、超时的惩罚幅度递增），使反馈同时携带吞吐收益与排队/丢包代价信息。

\textbf{相对性（Relativity）}：由于逐事件反馈几乎恒为正，绝对幅值不携带调节信息；Swift只使用反馈相对其自身慢速基线的\emph{变化量}驱动$\alpha$微调（详见第\ref{sec:adaptive-alpha}节），从而避免固定绝对阈值导致的单向漂移，并使同一组超参数可以跨带宽与RTT差异极大的路径复用。

\subsection{规则策略表示}

Swift的决策策略采用\emph{确定性规则}表示，而非随机策略或神经网络策略。这种选择的考量如下：

\textbf{确定性}：给定相同的状态向量，决策输出完全确定，行为可逐分支解释、可离线推演，实验可精确复现。

\textbf{参数化 vs.\ 非参数化}：神经网络策略是参数化策略，表示能力强但依赖训练与推理设施；基于规则的策略计算开销为常数级、无推理依赖，表示能力受人工规则结构限制。Swift选择后者，以可解释性换取工程可控性，其策略表示能力的局限与可能的增强路线见第\ref{Chap:chap5}章。

\textbf{离线设计 + 在线参数自适应}：规则主体离线设计，不经过在线学习；控制参数（乘性增加因子$\alpha$等）在运行期由反馈与RTT信号小步在线调整，使算法在保持确定性主路径的同时获得对环境变化的适应能力。

\section{11维观测空间详细设计}

本节详细描述Swift算法的11维观测空间设计，这是Swift区别于现有方案的核心创新之一。

\textbf{记法约定}：为与ns-3侧OpenGym实现（\path{contrib/opengym/examples/swift-tcp/tcp-swift-env.cc}中\texttt{parameterNum = 15}）保持一致，每次回调实际通过跨进程传输容器传递15个元素，其中前4个（socketUuid、envType、simTime、nodeId）为跨进程路由用的元数据通道，不参与AIMD/ECN响应/$\alpha$自适应等窗口决策；后11个（ssthresh、cwnd、segSize、segmentsAcked、bytesInFlight、lastRtt、minRtt、callbackType、caState、caEvent、ecnState）构成参与决策的有效状态子空间。故下文所称"11维观测空间"特指该有效子空间；表~\ref{tab:observation}同时列出元数据通道以便审阅者核对完整布局。

\begin{table}[htbp]
\centering
\caption{Swift 11维观测空间参数定义}
\label{tab:observation}
\begin{tabular}{|c|l|l|l|}
\hline
\textbf{索引} & \textbf{参数名称} & \textbf{数据类型} & \textbf{说明} \\
\hline
0 & socketUuid & uint64 & 套接字唯一标识符 \\
1 & envType & uint64 & 环境类型（0=事件驱动，1=时间驱动） \\
2 & simTime & uint64 & 仿真时间（微秒） \\
3 & nodeId & uint64 & 节点标识符 \\
4 & ssthresh & uint64 & 慢启动阈值（字节） \\
5 & cwnd & uint64 & 拥塞窗口大小（字节） \\
6 & segSize & uint64 & TCP报文段大小（字节） \\
7 & segmentsAcked & uint64 & 本次确认的报文段数量 \\
8 & bytesInFlight & uint64 & 在途字节数 \\
9 & lastRtt & uint64 & 最近往返时延（微秒） \\
10 & minRtt & uint64 & 最小往返时延（微秒） \\
11 & callbackType & uint64 & 回调函数类型 \\
12 & caState & uint64 & 拥塞算法状态 \\
13 & caEvent & uint64 & 拥塞算法事件 \\
14 & ecnState & uint64 & ECN状态 \\
\hline
\end{tabular}
\end{table}

\subsection{连接标识参数（索引0、3）}

\textbf{socketUuid（索引0）}是套接字的唯一标识符，用于区分不同的TCP连接。在多流并发场景下，每个流需要维护独立的状态信息，socketUuid是实现每流状态隔离的关键。socketUuid在连接建立时由系统分配，在连接生命周期内保持不变。

\textbf{nodeId（索引3）}是网络节点的标识符，用于区分不同的发送端节点。在多节点仿真场景下，不同节点可能运行不同的拥塞控制算法或具有不同的网络条件，nodeId用于识别当前决策所属的节点。

\subsection{时间信息参数（索引1、2）}

\textbf{envType（索引1）}是环境类型标识符。当前事件驱动实现中该字段固定为0，用于标识状态采集来自TCP回调事件路径。

\textbf{simTime（索引2）}是当前仿真时间，单位为微秒。仿真时间用于标识观测发生的事件顺序，并为跨进程观测和动作交互提供时间上下文。

\subsection{窗口状态参数（索引4、5）}

\textbf{ssthresh（索引4）}是慢启动阈值（Slow Start Threshold），单位为字节。ssthresh是TCP拥塞控制的核心状态变量之一，用于区分慢启动阶段和拥塞避免阶段：当$cwnd < ssthresh$时处于慢启动阶段，窗口指数增长；当$cwnd \geq ssthresh$时进入拥塞避免阶段，窗口线性增长。

\textbf{cwnd（索引5）}是拥塞窗口（Congestion Window）大小，单位为字节。cwnd是TCP拥塞控制的另一核心状态变量，决定了发送端在未收到确认的情况下最多可以发送的数据量。cwnd的调整是拥塞控制的核心任务。

\subsection{传输性能参数（索引6、7、8）}

\textbf{segSize（索引6）}是TCP报文段大小（Maximum Segment Size，MSS），单位为字节。本文仿真器按\texttt{mtu\_bytes} $-\;20\;-$ (IPv4首部 $+$ TCP首部)计算应用层数据单元，在MTU为1500字节、IPv4与TCP首部各20字节的默认配置下得到1440字节，其中额外扣除的20字节为仿真器为TCP选项字段预留的余量。segSize用于将以字节为单位的窗口度量换算为以报文段为单位的度量，并且是慢启动增量、加性递增步长以及窗口上下界的共同基准单位，因此该值的准确性直接决定全部窗口计算的量纲一致性。

\textbf{segmentsAcked（索引7）}是本次确认的报文段数量。在收到累积确认（Cumulative ACK）或选择性确认（Selective ACK，SACK）时，segmentsAcked表示新确认的报文段数量。segmentsAcked用于计算窗口增长量，例如在慢启动阶段，窗口增长量为$segmentsAcked \times MSS$。

\textbf{bytesInFlight（索引8）}是在途字节数（Bytes in Flight），即已发送但尚未确认的数据量。bytesInFlight用于计算管道利用率（Pipeline Utilization）：

\begin{equation}
utilization = \frac{bytesInFlight}{cwnd}
\end{equation}

当utilization较低时，表明发送端未能充分利用拥塞窗口，可能是因为应用层数据不足或接收端通告窗口限制。

\subsection{延迟测量参数（索引9、10）}

\textbf{lastRtt（索引9）}是最近测量的往返时延（Round-Trip Time），单位为微秒。RTT是端到端延迟的直接度量，包括传播延迟（Propagation Delay）、排队延迟（Queuing Delay）、处理延迟（Processing Delay）和传输延迟（Transmission Delay）四个组成部分。

\textbf{minRtt（索引10）}是历史最小往返时延，单位为微秒。minRtt作为路径传播时延的保守下界，用于估计带宽延迟积（BDP）和计算RTT膨胀程度。RTT膨胀比定义为：

\begin{equation}
inflation\_ratio = \frac{lastRtt}{minRtt}
\end{equation}

当inflation\_ratio接近1时，表明网络排队延迟较小；当inflation\_ratio较大时，表明网络存在显著的排队延迟。

\subsection{回调类型参数（索引11）}

\textbf{callbackType（索引11）}标识当前回调的类型，用于区分不同类型的事件。定义如下：

\begin{itemize}
    \item 值为0：GetSsThresh回调，表示协议栈请求更新慢启动阈值，触发原因需结合caState与ecnState区分
    \item 值为1：IncreaseWindow回调，表示收到正常确认并进入窗口调整路径
\end{itemize}

callbackType是多信号融合拥塞检测的关键输入之一。当callbackType为0时，可以确认协议栈正在执行窗口缩减流程，但不能仅凭该字段把原因判定为普通丢包。

\subsection{拥塞控制状态参数（索引12、13、14）}

\textbf{caState（索引12）}是拥塞算法状态（Congestion Algorithm State），反映TCP拥塞控制状态机的当前状态。Linux TCP实现定义了以下状态：

\begin{itemize}
    \item CA\_OPEN（值为0）：开放状态，正常传输，无拥塞迹象
    \item CA\_DISORDER（值为1）：无序状态，收到重复确认或乱序确认
    \item CA\_CWR（值为2）：拥塞窗口缩减状态，响应ECN标记
    \item CA\_RECOVERY（值为3）：恢复状态，快速恢复过程中
    \item CA\_LOSS（值为4）：丢失状态，检测到超时丢包
\end{itemize}

\textbf{caEvent（索引13）}是拥塞算法事件（Congestion Algorithm Event），表示最近发生的拥塞相关事件。Linux TCP实现定义了以下事件：

\begin{itemize}
    \item CA\_EVENT\_TX\_START（值为0）：传输开始事件
    \item CA\_EVENT\_CWND\_RESTART（值为1）：拥塞窗口重启事件
    \item CA\_EVENT\_COMPLETE\_CWR（值为2）：CWR完成事件
    \item CA\_EVENT\_LOSS（值为3）：丢失事件
    \item CA\_EVENT\_ECN\_NO\_CE（值为4）：无CE标记的ECN事件
    \item CA\_EVENT\_ECN\_IS\_CE（值为5）：有CE标记的ECN事件
\end{itemize}

\textbf{ecnState（索引14）}是显式拥塞通知状态（ECN State），反映ECN协商和标记的状态。定义如下：

\begin{itemize}
    \item ECN\_DISABLED（值为0）：ECN禁用
    \item ECN\_IDLE（值为1）：ECN空闲
    \item ECN\_CE\_RCVD（值为2）：收到CE标记
    \item ECN\_SENDING\_ECE（值为3）：正在发送ECE标志
    \item ECN\_ECE\_RCVD（值为4）：收到ECE标志
    \item ECN\_CWR\_SENT（值为5）：已发送CWR标志
\end{itemize}

caState、caEvent与ecnState为拥塞类型判定提供协议栈语义。Swift将这些字段与窗口、在途字节和RTT观测结合，用于区分超时、ECN驱动的窗口缩减与普通丢包；该观测集合面向当前控制规则设计，不代表覆盖TCP协议栈的全部状态。

\section{多信号融合拥塞检测算法}\label{sec:multi-signal}

Swift算法的核心创新之一是多信号融合拥塞检测方法。该方法综合利用丢包回调、ECN状态和TCP拥塞状态机信息，采用显式信号优先级判断拥塞类型，并将RTT膨胀保留给参数自适应调节，而不把RTT膨胀作为独立降窗触发条件。

\subsection{拥塞信号分类}

Swift将拥塞信号分为四个类别，每一类信号在"发生时刻—可靠性—语义粒度"三个维度上具有不同特征，这构成了后续差异化响应的判据基础。

\textbf{（1）显式降窗信号}：当callbackType为0（GetSsThresh回调）时，表明ns-3的TCP协议栈已经自行判定需要缩减发送窗口并向拥塞控制算法索取新的慢启动阈值。该回调是协议栈内部唯一的"权威降窗入口"，其触发条件包含快速重传检测到的丢包、重传超时以及ECN回显驱动的CWR进入三种情形，因此不能笼统地等同于"丢包"。Swift在该回调内部再依据caState和ecnState做二次判别，这是本文与既有方案在信号解释粒度上的关键差异。

\textbf{（2）ECN拥塞经历信号}：当ecnState取值为ECN\_CE\_RCVD（值为2，本端收到被标记CE的数据包）或ECN\_ECE\_RCVD（值为4，本端收到对端回显的ECE标志）时，表明路径上的队列长度已经越过AQM的标记门限，但尚未溢出。ECN是唯一由网络设备主动生成、且不依赖丢包或RTT推断的拥塞证据，其时序显著早于丢包，因此适合作为早期预警。Swift直接读取协议栈的ECN状态字段，而不维护标记事件的时间戳队列与标记率统计，从而避免了标记率阈值这一额外的敏感参数。

\textbf{（3）拥塞控制状态机信号}：caState刻画了协议栈当前所处的恢复阶段。其中CA\_LOSS（值为4）指示重传超时，是严重程度最高的信号；CA\_CWR（值为2）指示协议栈正因ECN回显而缩减窗口，在语义上属于ECN响应而非丢包响应；CA\_DISORDER与CA\_RECOVERY则指示乱序或快速恢复过程，不构成独立的降窗触发条件。将caState纳入状态向量使决策模块能够在同一个GetSsThresh回调内区分"超时""ECN驱动的CWR"与"普通丢包"三条语义完全不同的路径。

\textbf{（4）延迟膨胀信号}：RTT膨胀比（lastRtt/minRtt）反映瓶颈队列的驻留程度。该信号连续可微、采样密度高，适合驱动参数的渐进调节；但它同时混杂了无线调度抖动、ACK压缩、路径切换和终端侧处理延迟等非拥塞因素，在远距离与异构终端场景中误判概率明显高于前三类信号。因此Swift不将其用作独立的乘性递减触发条件，而仅用于$\alpha$自适应与慢启动退出判定。

\subsection{分层优先级机制}

不同类型的拥塞信号具有不同的可靠性和紧迫程度，Swift采用分层优先级机制进行处理：

\textbf{第一优先级：显式降窗回调的三路语义判别}

\begin{algorithm}[htbp]
\caption{显式降窗回调的三路语义判别}
\begin{algorithmic}[1]
\Require 回调类型 $callbackType$，拥塞算法状态 $caState$，ECN状态 $ecnState$
\If{$callbackType = 0$} \Comment{GetSsThresh回调}
    \State $is\_congested \leftarrow True$
    \If{$caState = CA\_LOSS$}
        \State $congestion\_type \leftarrow TIMEOUT$ \Comment{重传超时}
    \ElsIf{$ecnState \in \{ECN\_CE\_RCVD, ECN\_ECE\_RCVD\}$ \textbf{or} $caState = CA\_CWR$}
        \State $congestion\_type \leftarrow ECN$ \Comment{ECN回显驱动的CWR}
    \Else
        \State $congestion\_type \leftarrow LOSS$ \Comment{快速重传丢包}
    \EndIf
    \State \Return $(is\_congested, congestion\_type)$
\EndIf
\end{algorithmic}
\end{algorithm}

GetSsThresh回调本身只说明"协议栈决定降窗"，并不说明降窗的原因。ns-3的\texttt{TcpSocketBase}在三种情形下都会进入该回调：快速重传确认丢包、重传定时器超时（此时caState被置为CA\_LOSS）、以及收到ECE回显后调用\texttt{EnterCwr}（此时caState被置为CA\_CWR）。若不加区分地把整个回调映射为"丢包"，则ECN提供的早期预警在响应环节会被当作队列已溢出处理，从而抵消ECN机制本应带来的时序优势——这正是本算法早期版本存在的语义缺陷。当前实现在回调内部按"超时优先、ECN次之、普通丢包兜底"的顺序完成三路判别：CA\_LOSS对应最严重的路径级异常；ecnState处于CE\_RCVD或ECE\_RCVD、或caState已进入CA\_CWR的情形归入ECN类；其余情形归入普通丢包类。三类判别结果分别驱动0.50、0.75、0.70三个窗口保留因子，使响应力度与信号严重程度保持单调一致。

\textbf{第二优先级：ECN拥塞经历信号检测}

\begin{algorithm}[htbp]
\caption{ECN拥塞经历信号检测}
\begin{algorithmic}[1]
\Require ECN状态 $ecnState$
\If{$ecnState \in \{ECN\_CE\_RCVD, ECN\_ECE\_RCVD\}$}
    \State $is\_congested \leftarrow True$
    \State $congestion\_type \leftarrow ECN$
    \State \Return $(is\_congested, congestion\_type)$
\EndIf
\end{algorithmic}
\end{algorithm}

当ecnState为ECN\_CE\_RCVD（收到CE标记）或ECN\_ECE\_RCVD（收到ECE回显）时，判定发生ECN拥塞。ECN拥塞采用早期预警响应策略，当前实现的保留因子为0.75；该值比普通丢包响应更温和，同时能够更及时释放队列压力。

\subsection{信号优先级与误判抑制策略}

Swift的拥塞检测采用显式信号优先级，而不是将所有协议栈状态都作为独立拥塞触发条件。当前实现中，显式丢包回调具有最高优先级；ECN CE/ECE状态是第二优先级的早期拥塞信号；RTT膨胀、CA\_CWR和CA\_RECOVERY等状态不直接触发降窗，而用于解释网络状态或影响后续参数调节。

\textbf{（1）显式丢包优先}：当回调类型为GetSsThresh时，算法认为TCP协议栈已经进入丢包相关处理路径，并根据CA状态区分普通丢包和超时丢包。

\textbf{（2）ECN直接触发早期响应}：当ECN状态为ECN\_CE\_RCVD或ECN\_ECE\_RCVD时，算法将其视为有效早期拥塞信号并触发ECN响应，同时累加ECN计数器用于后续状态分析。

\textbf{（3）RTT膨胀不单独触发降窗}：RTT膨胀容易受到无线调度、ACK压缩和路径抖动影响。当前实现将其用于$\alpha$自适应调节和慢启动退出判断，而不是作为独立拥塞事件直接执行乘性递减。

\textbf{（4）降窗后冻结抑制反弹}：每次显式拥塞响应后，算法冻结后续若干ACK事件中的窗口增长，避免刚降窗后立即增窗造成二次排队。

\subsection{完整拥塞检测算法}

综合上述机制，Swift的完整多信号融合拥塞检测算法如下：

\begin{algorithm}[htbp]
\caption{多信号融合拥塞检测算法}
\begin{algorithmic}[1]
\Require 观测状态 $obs$，流状态 $state$
\Ensure 是否拥塞 $is\_congested$，拥塞类型 $type$
\State \Comment{第一优先级：显式降窗回调的三路语义判别}
\If{$obs.callbackType = 0$} \Comment{GetSsThresh回调}
    \If{$obs.caState = CA\_LOSS$}
        \State $state.loss\_count \leftarrow state.loss\_count + 1$
        \State \Return $(True, TIMEOUT)$ \Comment{重传超时}
    \EndIf
    \If{$obs.ecnState \in \{ECN\_CE\_RCVD, ECN\_ECE\_RCVD\}$ \textbf{or} $obs.caState = CA\_CWR$}
        \State $state.ecn\_count \leftarrow state.ecn\_count + 1$
        \State \Return $(True, ECN)$ \Comment{ECN回显驱动的CWR}
    \EndIf
    \State $state.loss\_count \leftarrow state.loss\_count + 1$
    \State \Return $(True, LOSS)$ \Comment{快速重传丢包}
\EndIf
\State \Comment{第二优先级：IncreaseWindow路径上的ECN状态巡检}
\If{$obs.ecnState \in \{ECN\_CE\_RCVD, ECN\_ECE\_RCVD\}$}
    \State $state.ecn\_count \leftarrow state.ecn\_count + 1$
    \State \Return $(True, ECN)$
\EndIf
\State \Comment{未检测到拥塞}
\State \Return $(False, NONE)$
\end{algorithmic}
\end{algorithm}

该算法的时间复杂度为$O(1)$，仅涉及常数次整型比较，不引入任何队列扫描或统计窗口维护开销，因此可以在每个ACK事件上无条件执行。算法的两级结构也决定了拥塞类型判定的互斥性：一次回调最多产生一个拥塞类型，不存在多个信号同时触发导致窗口被重复缩减的情形。

\section{差异化拥塞响应机制}

Swift算法针对不同类型的拥塞信号采用差异化的响应策略，核心是差异化窗口保留因子（Window Retention Factor）机制。

\subsection{窗口保留因子设计}

窗口保留因子$\rho$定义为拥塞响应后新窗口与原窗口的比值：

\begin{equation}
new\_cwnd = \rho \cdot cwnd
\end{equation}

传统算法（如TCP Reno、CUBIC）对所有拥塞信号采用统一的窗口保留因子（通常为0.5，即窗口减半）。这种"一刀切"的策略无法区分不同类型和程度的拥塞，在轻度拥塞时响应过度，在严重拥塞时响应不足。

Swift根据拥塞类型设计了差异化的窗口保留因子：

\begin{table}[htbp]
\centering
\caption{差异化窗口保留因子}
\label{tab:retention}
\begin{tabular}{|l|c|l|}
\hline
\textbf{拥塞类型} & \textbf{保留因子} & \textbf{设计理由} \\
\hline
丢包（Loss） & 0.70 & 丢包是严重拥塞信号，需要较大幅度降低窗口 \\
ECN标记（ECN） & 0.75 & ECN是早期信号，提前释放队列压力 \\超时丢失（Timeout） & 0.50 & 超时表示最严重拥塞，需最大幅度窗口缩减 \\
\hline
\end{tabular}
\end{table}

\subsection{差异化响应的理论依据}

差异化窗口保留因子的设计基于以下理论分析：

\textbf{（1）拥塞信号的时序特性}

ECN标记是最早的拥塞信号，通常在队列长度超过低阈值时触发，此时网络可能只是轻度拥塞。丢包是较晚的拥塞信号，通常在队列溢出时发生，表明网络已经严重拥塞。超时丢包是最晚也最严重的拥塞信号，通常意味着多个数据包连续丢失或网络路径严重中断，需要最大力度的窗口缩减。

根据时序关系和严重程度，Swift设计了三级差异化响应：ECN标记采用早期预警响应（保留因子0.75），丢包采用较激进的响应（保留因子0.70），超时采用最激进的响应（保留因子0.50）。这种阶梯式的响应力度设计用于匹配不同严重程度的拥塞信号，同时通过降窗后冻结机制抑制窗口立即反弹。

\textbf{（2）信号可靠性分析}

丢包通常表示报文未被成功交付，但在无线或切换场景中不一定由队列拥塞引起；ECN标记的语义则取决于网络设备是否启用标记及其阈值配置。因此，当前实现对普通丢包采用0.70的保留因子，对ECN采用0.75的保留因子，并把超时作为更严重的恢复状态处理。这些取值属于人工设定，其独立效果仍需消融实验检验。

\textbf{（3）吞吐量-延迟权衡}

较低的窗口保留因子能够更快地缓解拥塞、降低延迟，但可能降低后续吞吐量。较高的窗口保留因子有利于维持在途数据量，但可能延长拥塞持续时间。Swift的差异化设计用于在这两类效应之间进行折中，其具体效果由第~\ref{Chap:chap4}章的整体实验结果讨论，不作最优性断言。

\subsection{慢启动阈值更新}

除了拥塞窗口，拥塞响应还需要更新慢启动阈值。当前实现首先以$\rho \times cwnd$计算新的拥塞窗口，并设置最小窗口边界$W_{min}=\max(4\times MSS,1)$。慢启动阈值采用BDP锚定但不迁就队列膨胀的策略：

\begin{equation}
ref = \min(cwnd, \max(BDP, W_{min}))
\end{equation}

\begin{equation}
new\_ssthresh = \max(\rho \times ref, W_{min}, new\_cwnd)
\end{equation}

该设计避免在队列已膨胀时把慢启动阈值设置到过高水平，同时保证阈值不低于新的拥塞窗口，防止拥塞恢复阶段出现窗口停滞。

此外，Swift引入连续缩减保护和降窗后冻结两项稳定性机制，二者作用于不同的时间尺度且相互配合。

\textbf{连续缩减保护}针对的是"窗口死亡螺旋"：在突发拥塞或无线批量误码期间，协议栈可能在很短时间内连续触发多次降窗回调，若每次都执行乘性递减，窗口将以指数速度塌缩至下界，随后需要数十个RTT才能重新填满管道。算法为每条流维护连续缩减计数器$d$，每次执行拥塞响应时$d$自增；当$d$超过$D_{max} = 3$时，响应退化为"保持当前窗口"而不再继续缩减。取$D_{max} = 3$的依据是：连续三次丢包响应后窗口仅余$0.70^3 \approx 34.3\%$，已经低于传统算法单次窗口减半的结果，继续缩减的边际收益迅速递减而恢复代价持续上升。

\textbf{降窗后冻结}针对的是"降窗即反弹"：乘性递减只减小了发送窗口，瓶颈队列中已经排队的数据仍需一段时间才能排空；若在此期间立即恢复窗口增长，新注入的数据会与残留队列叠加，导致排队延迟不降反升，并可能诱发第二次拥塞事件。因此任一显式降窗后，算法冻结后续4个ACK事件中的窗口增长，为队列排空留出时间窗口。

两项机制的交互存在一个容易被忽视的细节：连续缩减计数器$d$的清零时机。当前实现只在"窗口确实发生增长"时才把$d$置零，即清零动作位于冻结判断\emph{之后}；若把清零动作前置到冻结判断之前，则冻结期内每个被冻结的ACK事件都会误将$d$清零，$D_{max}$计数被反复复位，连续缩减保护实际失效——这恰是本算法早期版本存在的语义缺陷。将清零动作与真实的窗口增长绑定，保证了"只有当算法确认网络已经恢复到可以增窗的状态时，才认为连续拥塞序列已经结束"。

\subsection{差异化拥塞响应算法}

\begin{algorithm}[htbp]
\caption{差异化拥塞响应算法}
\begin{algorithmic}[1]
\Require 当前窗口 $cwnd$，当前阈值 $ssthresh$，拥塞类型 $type$，连续缩减计数 $consec\_dec$
\Ensure 新窗口 $new\_cwnd$，新阈值 $new\_ssthresh$
\State $consec\_dec \leftarrow consec\_dec + 1$
\State \Comment{连续缩减保护：超过3次后保持窗口}
\If{$consec\_dec > 3$}
    \State \Return $(cwnd, cwnd)$ \Comment{保持当前窗口不再缩减}
\EndIf
\State \Comment{根据拥塞类型选择窗口保留因子}
\If{$type = TIMEOUT$}
    \State $\rho \leftarrow 0.50$ \Comment{超时：最严厉响应}
\ElsIf{$type = LOSS$}
    \State $\rho \leftarrow 0.70$ \Comment{丢包：较激进响应}
\ElsIf{$type = ECN$}
    \State $\rho \leftarrow 0.75$ \Comment{ECN：早期预警响应}
\Else
    \State $\rho \leftarrow 1.0$ \Comment{无拥塞，不调整}
\EndIf
\State \Comment{计算新窗口}
\State $new\_cwnd \leftarrow \lfloor \rho \times cwnd \rfloor$
\State \Comment{应用窗口下限约束}
\State $new\_cwnd \leftarrow \max(new\_cwnd, 4 \times MSS)$
\State \Comment{计算新慢启动阈值：以BDP和当前窗口的较小有效参考为锚定}
\State $ref \leftarrow \min(cwnd, \max(bdp, 4 \times MSS))$
\State $new\_ssthresh \leftarrow \max(\lfloor \rho \times ref \rfloor, new\_cwnd)$
\State \Return $(new\_cwnd, new\_ssthresh)$
\end{algorithmic}
\end{algorithm}

\section{启发式反馈自适应的参数调优}
\label{sec:adaptive-alpha}

Swift算法实现了启发式反馈驱动的参数自适应机制，根据多种因素动态调整乘性增加因子（Multiplicative Increase Factor，$\alpha$）。

\subsection{乘性增加因子的作用}

在Swift的融合控制策略中，乘性增加因子$\alpha$控制拥塞避免阶段的窗口增长速度：

\begin{equation}
target\_ref = \alpha \times BDP + \gamma \times MSS
\end{equation}

$\alpha$的取值决定了窗口围绕BDP目标的激进程度：较大的$\alpha$使目标窗口更接近或略高于估计BDP，能够更快地利用可用带宽，但也更容易形成排队；较小的$\alpha$使目标窗口低于或接近BDP，有利于排队释放和延迟恢复。

\subsection{多因素调整机制}

Swift根据以下因素动态调整$\alpha$：

\textbf{（1）基于RTT膨胀比的路径感知调整}

RTT膨胀比反映网络排队延迟程度：

\begin{equation}
inflation\_ratio = \frac{lastRtt}{minRtt}
\end{equation}

调整规则采用与最小RTT相关的动态阈值。首先根据最小RTT计算松弛量：
\begin{equation}
s = 0.3 + 0.15\sqrt{\max(0, RTT_{min,ms}-1)}
\end{equation}
并构造$1+s$、$1+2s$和$1+4s$三个路径感知阈值。短RTT路径对排队更敏感，阈值更紧；长RTT路径允许更大的正常抖动，阈值更宽。具体调整规则为：
\begin{itemize}
    \item 当$inflation\_ratio < 1+s$时，认为排队较小，增大$\alpha$；若最小RTT超过5 ms，则增长步长为0.02，否则为0.03。
    \item 当$1+s \leq inflation\_ratio < 1+2s$时，认为排队适中，温和增大$\alpha$：$\alpha \leftarrow \alpha + 0.01$。
    \item 当$inflation\_ratio > 1+4s$时，认为排队较重，减小$\alpha$；若最小RTT超过5 ms，则减小步长为0.03，否则为0.02。
\end{itemize}

这种路径感知响应机制能够避免固定阈值在远距离路径上过早收紧窗口，同时保留低RTT路径中对排队增长的敏感性。

\textbf{（2）基于性能反馈值基线相对趋势的调整}

Swift利用环境合成的性能反馈值评估当前控制策略的相对优劣。反馈值在每个正常确认事件上由确认段数带来的吞吐收益减去RTT膨胀惩罚构成，因此在绝大多数ACK事件上取正值。这一分布特征决定了不能用固定的绝对阈值判断"反馈是否变好"：若采用形如$R_{ema} > 0.5$的判据，则该条件在几乎每个ACK事件上都成立，$\alpha$会被持续推向上界$\alpha_{max}$，自适应机制退化为单向棘轮，失去调节意义。

当前实现改为\textbf{基线相对判据}，即同时维护快、慢两条指数移动平均，用快线相对慢线的偏离量而非绝对水平来判断趋势：

\begin{equation}
R_{fast} \leftarrow (1 - \eta_f) \times R_{fast} + \eta_f \times reward, \quad \eta_f = 0.15
\end{equation}

\begin{equation}
R_{base} \leftarrow (1 - \eta_b) \times R_{base} + \eta_b \times reward, \quad \eta_b = 0.02
\end{equation}

其中$R_{fast}$为快线，反映最近若干次决策的即时效果；$R_{base}$为慢线基线，反映当前网络条件下的长期反馈水平。两条EMA在收到首个反馈值时用同一数值初始化，避免冷启动阶段的虚假偏离。判定容差按基线幅值自适应缩放：

\begin{equation}
m = \max(0.25,\; 0.1 \times |R_{base}|)
\end{equation}

调整规则为：
\begin{itemize}
    \item 当$R_{fast} > R_{base} + m$时，表明近期决策相对该路径的常态水平确有改善，适度提高激进度：$\alpha \leftarrow \alpha + 0.01$。
    \item 当$R_{fast} < R_{base} - 4m$时，表明近期决策显著劣于常态水平，回退激进度：$\alpha \leftarrow \alpha - 0.01$。
\end{itemize}

下调阈值取$4m$而上调阈值取$m$，构成有意的非对称设计：反馈序列的负向尖峰主要由单次丢包或超时惩罚引起，而这类事件已经由差异化拥塞响应机制单独处理，若在$\alpha$层面再次同步收缩会造成响应叠加；采用较宽的下调死区可以把$\alpha$的收缩限定在持续性劣化而非瞬时惩罚上。基线相对判据的另一优点是与路径无关：不同带宽和RTT的路径其反馈绝对幅值相差可达数个数量级，而相对偏离量是无量纲的，因此同一组超参数可以直接跨场景复用，无需按链路类型重新标定。

\textbf{（3）基于连续增长趋势的调整}

当窗口连续增长超过8次时，表明网络状况持续良好，给予稳定增长的正向反馈：
\begin{equation}
\alpha \leftarrow \alpha + 0.01 \quad \text{如果连续增长次数} > 8
\end{equation}

\subsection{参数边界约束}

为保证控制策略的稳定性，$\alpha$的取值被限制在预设范围内：

\begin{equation}
\alpha \in [\alpha_{min}, \alpha_{max}] = [0.85, 1.30]
\end{equation}

$\alpha$的默认初始值为1.10。下界0.85为无线、蜂窝或浅缓冲路径提供更强的排队释放空间；上界1.30限制窗口相对于BDP的长期扩张幅度，避免慢性队列积累。

\subsection{自适应参数调优算法}

\begin{algorithm}[htbp]
\caption{自适应参数调优算法}
\begin{algorithmic}[1]
\Require 当前$\alpha$，观测状态$obs$，性能反馈值$feedback$，历史状态$history$
\Ensure 更新后的$\alpha$
\State \Comment{因素1：基于RTT膨胀比的路径感知调整}
\State $inflation\_ratio \leftarrow obs.lastRtt / obs.minRtt$
\State $min\_rtt\_ms \leftarrow obs.minRtt / 1000$
\State $slack \leftarrow 0.3 + 0.15\sqrt{\max(0, min\_rtt\_ms-1)}$
\State $low \leftarrow 1 + slack$, $mid \leftarrow 1 + 2\times slack$, $high \leftarrow 1 + 4\times slack$
\If{$inflation\_ratio < low$}
    \State $ramp \leftarrow 0.02$ if $min\_rtt\_ms > 5$ else $0.03$
    \State $\alpha \leftarrow \alpha + ramp$
    \State $history.consecutive\_growth \leftarrow history.consecutive\_growth + 1$
\ElsIf{$inflation\_ratio < mid$}
    \State $\alpha \leftarrow \alpha + 0.01$
\ElsIf{$inflation\_ratio > high$}
    \State $step \leftarrow 0.03$ if $min\_rtt\_ms > 5$ else $0.02$
    \State $\alpha \leftarrow \alpha - step$
    \State $history.consecutive\_growth \leftarrow 0$
\EndIf
\State \Comment{因素2：基于反馈值基线相对偏离的趋势调整}
\If{首次收到反馈值}
    \State $R_{fast} \leftarrow reward$, $R_{base} \leftarrow reward$
\Else
    \State $R_{fast} \leftarrow 0.85 \times R_{fast} + 0.15 \times reward$
    \State $R_{base} \leftarrow 0.98 \times R_{base} + 0.02 \times reward$
\EndIf
\State $m \leftarrow \max(0.25, 0.1 \times |R_{base}|)$
\If{$R_{fast} > R_{base} + m$}
    \State $\alpha \leftarrow \alpha + 0.01$
\ElsIf{$R_{fast} < R_{base} - 4 \times m$}
    \State $\alpha \leftarrow \alpha - 0.01$
\EndIf
\State \Comment{因素3：基于连续稳定增长的反馈}
\If{$history.consecutive\_growth > 8$}
    \State $\alpha \leftarrow \alpha + 0.01$
\EndIf
\State \Comment{应用边界约束}
\State $\alpha \leftarrow \max(\min(\alpha, 1.30), 0.85)$
\State \Return $\alpha$
\end{algorithmic}
\end{algorithm}

\section{融合控制策略}

Swift采用融合控制策略计算目标拥塞窗口，将基于带宽延迟积（BDP）估计的速率控制与基于当前窗口的保守控制相结合。

\subsection{带宽延迟积估计}

带宽延迟积（Bandwidth-Delay Product，BDP）是网络链路能够容纳的最大在途数据量，是拥塞控制的重要参考指标：

\begin{equation}
BDP = Bandwidth \times RTT
\end{equation}

Swift采用"时间窗口投递速率采样 + 窗口化最大值过滤"两级结构估计瓶颈带宽，再与最小RTT相乘得到BDP：

\begin{equation}
BDP = \max_{i \in W_{bw}}(DR_i) \times RTT_{min}
\label{eq:bdp}
\end{equation}

其中$W_{bw}$为容量40个样本的滑动窗口，$RTT_{min}$为该流历史最小RTT（同时接受协议栈上报的minRtt与本地观测的lastRtt作为候选下界）。

\textbf{时间窗口投递速率采样}。单个投递速率样本由\emph{累计确认字节数在一段时间窗口上的增量}给出。算法为每条流维护累计确认字节数$B(t)$，并在队列中记录二元组$(t, B(t))$；每次收到确认后先更新$B(t)$，再淘汰窗口左端超出时间跨度的样本，最后按下式计算样本：

\begin{equation}
DR = \frac{B(t) - B(t - \Delta)}{\Delta}, \quad \Delta = \mathrm{clamp}(2 \times RTT_{min},\; 5~\text{ms},\; 1~\text{s})
\label{eq:dr}
\end{equation}

窗口跨度$\Delta$取最小RTT的两倍，是为了让采样区间至少覆盖一个完整的ACK时钟周期：只有当窗口跨度不小于一个RTT时，窗口内的确认字节总量才等于瓶颈在该时段实际排空的数据量，$DR$才具有"瓶颈排空速率"的物理含义。下界5 ms用于抑制微秒级RTT路径上因窗口过短而放大的采样噪声，上界1 s用于保证带宽下降时估计能够及时跟随。

\textbf{为何不能按单次确认计算投递速率}。本算法早期版本采用形如$DR = segmentsAcked \times segSize / lastRtt$的单次确认公式。该式的分子只包含\emph{本次}确认的字节数，而分母却是\emph{整个}RTT，因此它实际估计的是"若一个RTT内只到达这一个确认时的速率"，与真实瓶颈速率相差约等于每RTT的确认个数$N_{ack} \approx cwnd/(segSize \times b)$（其中$b$为延迟确认因子）。在高BDP路径上$N_{ack}$可达数百，估计值因而被系统性压低同一量级，式~\eqref{eq:bdp}给出的BDP随之崩塌；此时拥塞避免阶段的目标窗口$\alpha \times BDP$长期低于当前窗口，窗口只能停留在安全下界$200 \times MSS$上，吞吐量被硬性锁定为$200 \times MSS / RTT$。这一失效模式在第~\ref{Chap:chap4}章的跨域WAN场景中留下了可验证的痕迹，也是本文将该估计器整体替换为式~\eqref{eq:dr}的直接原因。式~\eqref{eq:dr}的分子与分母取自同一时间区间，量纲自洽，不再依赖每RTT确认个数。

\textbf{窗口化最大值过滤}。将$DR$样本压入容量40的滑动窗口并取窗口内最大值，其作用有三：（1）投递速率样本天然是瓶颈能力的\emph{下界估计}——确认稀疏、应用层供数不足或跨流量抢占都只会使样本偏小而不会偏大，因此取最大值是对真实瓶颈的一致估计；（2）滑动窗口限制了历史记忆长度，使估计能够跟随带宽的真实下降，避免全局最大值带来的永久性高估；（3）窗口容量取40而非较小值，是为了让长RTT路径在窗口内容纳足够多的采样周期——GEO卫星链路单个RTT可达600 ms，过短的窗口会在瓶颈被充分采样之前就发生淘汰，导致估计值在低位振荡。

此外，流级吞吐量趋势由指数移动平均跟踪，用于状态观察与性能记录：

\begin{equation}
throughput\_ema \leftarrow 0.9 \times throughput\_ema + 0.1 \times DR
\end{equation}

BDP估计的完整流程如算法~\ref{alg:bdp-est}所示。当$\max(DR)$尚未建立或$RTT_{min}$无效时，BDP回退为当前拥塞窗口，使算法在连接建立初期退化为常规的窗口驱动行为而非产生无意义的目标值。

\begin{algorithm}[htbp]
\caption{时间窗口投递速率与BDP估计}
\label{alg:bdp-est}
\begin{algorithmic}[1]
\Require 流状态 $state$，观测状态 $obs$
\Ensure BDP估计值
\If{$obs.lastRtt > 0$}
    \State $state.min\_rtt \leftarrow \min(state.min\_rtt, obs.lastRtt)$
\EndIf
\If{$0 < obs.minRtt < state.min\_rtt$}
    \State $state.min\_rtt \leftarrow obs.minRtt$
\EndIf
\If{$obs.segmentsAcked > 0$ \textbf{and} $obs.segSize > 0$}
    \State $state.acked\_total \leftarrow state.acked\_total + obs.segmentsAcked \times obs.segSize$
    \State $state.samples$.Append($(obs.simTime, state.acked\_total)$)
    \State $\Delta \leftarrow \mathrm{clamp}(2 \times state.min\_rtt,\; 5~\text{ms},\; 1~\text{s})$
    \While{$|state.samples| \geq 2$ \textbf{and} $obs.simTime - state.samples[0].t > \Delta$}
        \State $state.samples$.PopFront()
    \EndWhile
    \State $span \leftarrow obs.simTime - state.samples[0].t$
    \State $delivered \leftarrow state.acked\_total - state.samples[0].B$
    \If{$|state.samples| \geq 2$ \textbf{and} $span > 0$ \textbf{and} $delivered > 0$}
        \State $DR \leftarrow delivered / span$
        \State $state.bw\_samples$.Append($DR$) \Comment{容量40的滑动窗口}
        \State $state.max\_bw \leftarrow \max(state.bw\_samples)$
    \EndIf
\EndIf
\If{$state.max\_bw > 0$ \textbf{and} $state.min\_rtt < \infty$}
    \State $state.bdp \leftarrow state.max\_bw \times state.min\_rtt$
\EndIf
\State \Return $state.bdp > 0$ ? $state.bdp$ : $\max(obs.cwnd, 1)$
\end{algorithmic}
\end{algorithm}

\subsection{慢启动阶段策略}

在慢启动阶段（$cwnd < ssthresh$），Swift采用指数增长策略以快速探测可用带宽：

\begin{equation}
target\_cwnd = 2 \times BDP
\end{equation}

\begin{equation}
increment = segmentsAcked \times MSS
\end{equation}

当窗口远低于BDP时，采用加速增长：

\begin{equation}
\text{如果~} cwnd < 0.3 \times BDP, \quad increment = 2 \times segmentsAcked \times MSS
\end{equation}

慢启动目标设为$2 \times BDP$而非更大值（如$3 \times BDP$），是为了避免慢启动阶段过度膨胀导致突发丢包。当窗口达到目标值后，自动退出慢启动进入拥塞避免阶段。

\subsection{拥塞避免阶段策略}

在拥塞避免阶段（$cwnd \geq ssthresh$），Swift采用稳健的窗口增长策略：

\begin{equation}
target\_ref = \alpha \times BDP + \gamma \times MSS
\end{equation}

其中，$\alpha$是乘性增加因子（通过自适应调优动态调整，基准值1.10），$\gamma = 1.0$是加性增加因子，表示每个ACK事件引入1个MSS的基础增量。拥塞避免阶段不再简单地以当前窗口作为下界持续取最大值，而是围绕$\alpha \times BDP$目标进行拉升或回落：当目标高于当前窗口时，以$\max(\gamma \times MSS, MSS, gap/16)$的步长向目标靠近；当目标低于当前窗口时，每个ACK事件按超出量的一半进行回落，以更快释放浅缓冲路径中的排队。

融合控制策略还考虑管道利用率因素。管道利用率定义为：

\begin{equation}
utilization = \frac{bytesInFlight}{cwnd}
\end{equation}

调整规则：
\begin{itemize}
    \item 仅当RTT未明显膨胀（$lastRtt < 1.3 \times minRtt$）时启用利用率补偿。
    \item 当$utilization < 0.4$时，短RTT路径额外增加$2 \times MSS$，长RTT路径额外增加$1 \times MSS$。
    \item 当$0.4 \leq utilization < 0.5$且路径不是长RTT路径时，额外增加$1 \times MSS$。
\end{itemize}

该设计只在管道确实未充分利用且队列未增长时加速，避免旧式利用率补偿在无线、蜂窝或浅缓冲路径中造成持续排队。

\subsection{窗口边界约束}

为保证窗口大小的合理性，Swift对计算结果应用边界约束：

\begin{equation}
cwnd \in [4 \times MSS, \max(4 \times BDP, 200 \times MSS)]
\end{equation}

下界$4 \times MSS$保证基本传输能力。上界设为BDP的4倍或200个MSS的较大值，用于在BDP估计暂时偏大或路径缓冲较浅时限制窗口长期膨胀。

\subsection{融合控制策略算法}

\begin{algorithm}[htbp]
\caption{融合控制策略算法}
\begin{algorithmic}[1]
\Require 观测状态$obs$，自适应参数$\alpha$，$\gamma = 1.0$，BDP估计$bdp$
\Ensure 新窗口$new\_cwnd$
\State \Comment{降窗后冻结：显式降窗后的若干ACK事件内保持窗口不增长}
\If{$state.freeze\_acks\_remaining > 0$}
    \State $state.freeze\_acks\_remaining \leftarrow state.freeze\_acks\_remaining - 1$
    \State \Return $obs.ssthresh, obs.cwnd$
\EndIf
\State \Comment{计算管道利用率}
\State $utilization \leftarrow obs.bytesInFlight / obs.cwnd$
\State \Comment{判断当前阶段}
\If{$obs.cwnd < obs.ssthresh$ \textbf{and} $in\_slow\_start$}
    \State \Comment{HyStart风格慢启动退出}
    \If{$obs.lastRtt > HystartThreshold(state) \times state.ss\_min\_rtt\_sample$}
        \State $state.in\_slow\_start \leftarrow False$
        \State \Return $\max(obs.ssthresh, obs.cwnd), obs.cwnd$
    \EndIf
    \State $target \leftarrow \max(2 \times bdp, 10 \times obs.segSize)$
    \State $increment \leftarrow obs.segmentsAcked \times obs.segSize$
    \If{$obs.cwnd < 0.3 \times bdp$ \textbf{and} $obs.lastRtt < 1.2 \times state.min\_rtt$}
        \State $increment \leftarrow 2 \times obs.segmentsAcked \times obs.segSize$
    \EndIf
    \State $new\_cwnd \leftarrow \min(obs.cwnd + increment, target)$
\Else
    \State \Comment{拥塞避免阶段：围绕$\alpha \times BDP$目标进行有界步长拉升或回落}
    \State $state.in\_slow\_start \leftarrow False$
    \State $target\_rate \leftarrow \lfloor \alpha \times bdp \rfloor$
    \State $gamma\_bytes \leftarrow \gamma \times obs.segSize$
    \If{$target\_rate > obs.cwnd$}
        \State $gap \leftarrow target\_rate - obs.cwnd$
        \State $step \leftarrow \max(gamma\_bytes, obs.segSize, gap / 16)$
        \State $new\_cwnd \leftarrow \min(obs.cwnd + step, target\_rate + gamma\_bytes)$
    \ElsIf{$target\_rate < obs.cwnd$}
        \State $excess \leftarrow obs.cwnd - target\_rate$
        \State $new\_cwnd \leftarrow \max(obs.cwnd - \max(excess / 2, obs.segSize), target\_rate)$
    \Else
        \State $new\_cwnd \leftarrow obs.cwnd + gamma\_bytes$
    \EndIf
    \State \Comment{仅在RTT未膨胀时启用管道利用率补偿；长RTT路径降低补偿强度}
    \If{$obs.lastRtt < 1.3 \times state.min\_rtt$ \textbf{and} $utilization < 0.4$}
        \State $new\_cwnd \leftarrow new\_cwnd + (state.min\_rtt > 5ms ? obs.segSize : 2 \times obs.segSize)$
    \ElsIf{$obs.lastRtt < 1.3 \times state.min\_rtt$ \textbf{and} $utilization < 0.5$ \textbf{and} $state.min\_rtt \leq 5ms$}
        \State $new\_cwnd \leftarrow new\_cwnd + obs.segSize$
    \EndIf
\EndIf
\State \Return $obs.ssthresh, new\_cwnd$
\end{algorithmic}
\end{algorithm}

\section{每流独立状态管理}

Swift设计了每流独立的状态管理机制，为每个TCP连接维护独立的状态数据结构，支持多流并发场景下的差异化控制。

\subsection{状态数据结构}

每个TCP连接维护以下状态数据，采用层次化组织方式：

\textbf{（1）历史度量缓冲区}：带宽估计维护两级队列。第一级是投递速率\emph{采样}队列，记录"仿真时刻—累计确认字节数"二元组，容量上限4096项，用于按时间窗口差分计算单个投递速率样本；每次更新时淘汰左端超出窗口跨度$\Delta = \mathrm{clamp}(2 \times RTT_{min}, 5~\text{ms}, 1~\text{s})$的样本，因此队列的实际长度随路径RTT与确认密度自适应变化。第二级是投递速率\emph{过滤}队列，容量40个样本，取窗口内最大值作为瓶颈带宽估计。RTT跟踪维护容量40个样本的最近RTT队列，同时维护该流的全局最小RTT作为传播延迟估计。BDP由最大带宽与最小RTT的乘积得到，并在估计尚未建立时回退为当前拥塞窗口。

\textbf{（2）拥塞事件计数数据}：状态结构维护丢包事件计数和ECN标记事件计数，用于记录每条流在仿真过程中的拥塞事件分布。

\textbf{（3）性能指标数据}：性能指标模块维护吞吐量的指数移动平均值（EMA，平滑因子0.1），用于跟踪流级别的性能变化趋势。

\textbf{（4）自适应参数数据}：自适应参数模块维护乘性增加因子$\alpha$（取值范围$[0.85, 1.30]$，默认值$1.10$）、加性增加因子$\gamma$（默认值$1.0$）、连续缩减计数器（用于连续缩减保护）、连续增长计数器（用于稳定增长反馈）、降窗后冻结ACK计数器（默认4个ACK事件）和慢启动阶段标记。

\subsection{状态查找与更新}

Swift使用socketUuid作为键，通过哈希表实现$O(1)$时间复杂度的状态查找和更新。哈希表的选择基于查找效率、动态扩展和内存效率三方面考量。

\begin{algorithm}[htbp]
\caption{流状态查找与更新}
\begin{algorithmic}[1]
\Function{LookupFlowState}{$manager$, $socket\_uuid$}
    \State $hash \leftarrow$ ComputeHash($socket\_uuid$)
    \State $bucket \leftarrow manager.buckets[hash \bmod manager.num\_buckets]$
    \For{each $entry$ in $bucket$}
        \If{$entry.key = socket\_uuid$}
            \State \Return $entry.value$
        \EndIf
    \EndFor
    \State \Return $null$ \Comment{未找到}
\EndFunction
\Function{UpdateFlowState}{$manager$, $socket\_uuid$, $obs$}
    \State $state \leftarrow$ GetOrCreateFlowState($manager$, $socket\_uuid$)
    \State $state.last\_update\_time \leftarrow obs[2]$
    \State UpdateMetrics($state.metrics\_buffer$, $obs$)
    \If{IsECNEvent($obs$)}
        \State $state.ecn\_tracker$.RecordEvent($obs[2]$)
    \EndIf
    \State UpdatePerformanceMetrics($state.performance$, $obs$)
\EndFunction
\end{algorithmic}
\end{algorithm}

\subsection{状态生命周期管理}

TCP连接的生命周期包括建立、传输和关闭三个阶段。Swift的状态管理遵循延迟初始化、增量更新、垃圾回收和状态持久化四个原则。

\begin{algorithm}[htbp]
\caption{流状态垃圾回收}
\begin{algorithmic}[1]
\Function{GarbageCollect}{$manager$, $current\_time$, $inactive\_threshold$}
    \State $to\_remove \leftarrow$ EmptyList()
    \For{each $(uuid, state)$ in $manager.flow\_states$}
        \State $idle\_time \leftarrow current\_time - state.last\_update\_time$
        \If{$idle\_time > inactive\_threshold$}
            \State $to\_remove$.Append($uuid$)
        \EndIf
    \EndFor
    \For{each $uuid$ in $to\_remove$}
        \State $manager.flow\_states$.Remove($uuid$)
        \State FreeResources($uuid$)
    \EndFor
    \State \Return $|to\_remove|$
\EndFunction
\end{algorithmic}
\end{algorithm}

\section{实现细节与优化技术}

本节详细介绍Swift算法的具体实现细节和优化技术。

\subsection{系统架构设计}

Swift算法的系统架构遵循模块化设计原则，采用分层结构实现各功能模块的解耦。整体架构分为仿真层、环境层、决策层和通信层四个层次。仿真层负责网络拓扑构建和数据包传输仿真；环境层作为协议栈与决策进程之间的桥梁；决策层实现Swift算法的核心逻辑，遵循状态隔离、分层决策和实时性原则；通信层负责仿真进程和决策进程之间的数据交换。

\begin{algorithm}[htbp]
\caption{Swift系统初始化流程}
\begin{algorithmic}[1]
\Require 网络拓扑配置 $topology$，仿真参数 $params$
\Ensure 初始化完成的Swift系统
\State $simulator \leftarrow$ CreateSimulator($params.duration$)
\State $network \leftarrow$ BuildTopology($topology$)
\State ConfigureQueues($network$, $params.queue\_size$)
\State $env \leftarrow$ CreateStateChannel()
\State RegisterCallbacks($env$, $network.tcp\_sockets$)
\State ConfigureObservationSpace($env$, $15$)
\State ConfigureActionSpace($env$, $2$) \Comment{ssthresh和cwnd}
\State $agent \leftarrow$ CreateSwiftAgent()
\State $agent.flow\_states \leftarrow$ EmptyHashMap()
\State $agent.\alpha_{default} \leftarrow 1.10$
\State $channel \leftarrow$ CreateMessageChannel()
\State ConnectEnvAndAgent($env$, $agent$, $channel$)
\State $simulator$.Run()
\end{algorithmic}
\end{algorithm}

\subsection{核心算法伪代码}

Swift决策模块的主控制循环负责协调各功能模块：处理来自环境层的状态采集结果，依据规则输出窗口决策，并把决策回写协议栈。

\begin{algorithm}[htbp]
\caption{Swift主控制循环}
\begin{algorithmic}[1]
\Require 环境接口 $env$，最大迭代次数 $max\_iterations$
\Ensure 控制结果序列
\State $iteration \leftarrow 0$
\State $flow\_manager \leftarrow$ InitFlowStateManager()
\While{$iteration < max\_iterations$ \textbf{and not} $env$.IsDone()}
    \State $obs \leftarrow env$.GetObservation()
    \State $socket\_uuid \leftarrow obs[0]$
    \State $flow\_state \leftarrow flow\_manager$.GetOrCreate($socket\_uuid$)
    \State UpdateFlowState($flow\_state$, $obs$)
    \State $(is\_congested, cong\_type, severity) \leftarrow$ DetectCongestion($obs$, $flow\_state$)
    \If{$is\_congested$}
        \State $action \leftarrow$ CongestionResponse($obs$, $cong\_type$, $flow\_state$)
    \Else
        \State UpdateAdaptiveParams($flow\_state$, $obs$)
        \State $action \leftarrow$ FusionControl($obs$, $flow\_state$)
    \EndIf
    \State $env$.ExecuteAction($action$)
    \State $iteration \leftarrow iteration + 1$
\EndWhile
\end{algorithmic}
\end{algorithm}

\subsection{数据结构设计}

Swift算法采用高效的数据结构支持其核心功能。\textbf{循环队列（Circular Buffer）}用于存储带宽与RTT历史样本，具有固定内存占用、$O(1)$时间复杂度的插入和删除操作、自动保持时间局部性等优势。\textbf{每流状态字典}以socketUuid为键隔离不同TCP连接的$\alpha$、BDP、拥塞事件计数、连续缩减计数和降窗后冻结计数，避免不同流之间的控制状态相互污染。

\subsection{性能优化技术}

为了提高Swift算法的执行效率，满足实时决策的要求，本文采用了多种性能优化技术：

\textbf{（1）延迟初始化策略}：流状态采用延迟初始化策略，只在首次访问时创建状态数据结构。

\begin{algorithm}[htbp]
\caption{延迟初始化状态获取}
\begin{algorithmic}[1]
\Function{GetOrCreateFlowState}{$manager$, $socket\_uuid$}
    \If{$socket\_uuid \notin manager.flow\_states$}
        \State $new\_state \leftarrow$ CreateEmptyFlowState()
        \State $new\_state.\alpha \leftarrow 1.10$
        \State $new\_state.min\_rtt \leftarrow \infty$
        \State $new\_state.max\_throughput \leftarrow 0$
        \State $new\_state.consecutive\_decreases \leftarrow 0$
        \State $new\_state.consecutive\_increases \leftarrow 0$
        \State $new\_state.in\_slow\_start \leftarrow True$
        \State $manager.flow\_states[socket\_uuid] \leftarrow new\_state$
    \EndIf
    \State \Return $manager.flow\_states[socket\_uuid]$
\EndFunction
\end{algorithmic}
\end{algorithm}

\textbf{（2）增量计算优化}：吞吐量、BDP等指标的计算采用增量方式，时间复杂度为$O(1)$。

\textbf{（3）分支路径清晰化}：拥塞检测、拥塞响应、慢启动和拥塞避免分别采用独立分支，便于将丢包、ECN、超时和普通ACK事件映射到可解释的窗口更新路径。

\textbf{（4）有界步长更新}：拥塞避免阶段的窗口拉升、回落和利用率补偿均采用有界步长，避免单个ACK事件造成窗口突变。

\textbf{（5）降窗后冻结}：显式降窗后冻结后续4个ACK事件中的窗口增长，减少窗口刚缩减后立即反弹造成的二次排队。

\subsection{状态—决策同步交互路径}
\label{sec:sync-path}

TcpSwift当前实验采用ZeroMQ同步请求—响应路径：ns-3仿真进程在TCP回调处生成状态向量（由ns3-gym/OpenGym框架的传输容器承载，仅作消息通道），Python决策进程返回慢启动阈值和拥塞窗口决策。该路径便于逐ACK解释窗口决策、复现实验数据并与\texttt{tcp\_swift.py}中的真实策略保持一致。由于该路径依赖同步消息交互，本文实验未将其与MTP并行仿真结合，也不报告MTP加速数据。

\subsection{正确性保证机制}
\label{sec:correctness}

Swift实现包含窗口边界、待应用动作时效性和异常观测回退等防护。这些机制限制当前控制器的输出范围并避免已知的陈旧动作问题，但不构成对任意跨进程输入或所有TCP状态组合的形式化正确性保证。

\textbf{（1）跨进程决策的时效性校验}。同步信道采用请求—应答模式，状态的采集时刻与决策的应用时刻之间存在一个消息往返间隔。在此间隔内协议栈的状态可能已经改变，最关键的情形是重传定时器超时。\texttt{tcp-swift-env.cc}在\texttt{CongestionStateSet}回调中检测到新状态为CA\_LOSS时，清除\texttt{m\_hasPendingCwnd}标志并丢弃尚未生效的窗口值，避免超时前的决策覆盖协议栈已更新的恢复状态。

\textbf{（2）动作边界约束}。Python控制器将拥塞窗口夹紧到$[4 \times MSS, \max(4 \times BDP, 200 \times MSS)]$区间。环境侧在接收动作后，以$2\times MSS$分别校验拥塞窗口和慢启动阈值的下界。该组合约束用于限制窗口过小或过度增长的风险；$200 \times MSS$项同时为BDP尚未建立的连接保留增长空间。

\textbf{（3）异常观测回退}。环境侧将未初始化或非正RTT编码为零，并在报文段大小为零时使用安全值；Python控制器在BDP尚未建立时以当前拥塞窗口作为回退值。

\textbf{（4）实验可复现性}。ns-3.40通过\texttt{RngRun}区分随机运行号，可对同一(场景, 协议)组合执行多次重复并汇总。第~\ref{Chap:chap4}章每个配置仅执行一次确定性运行，因此可按给定配置与运行号重复，但不能据此估计跨运行号方差。

\section{本章小结}

本章说明了Swift算法的系统结构及三项核心技术。

（1）\textbf{多信号拥塞状态感知}：ns-3环境通过15元素容器传输4项元数据与11项有效状态，并结合GetSsThresh、IncreaseWindow、PktsAcked、CongestionStateSet和CwndEvent回调更新状态。在GetSsThresh路径中，控制器依据caState与ecnState区分超时、ECN驱动的窗口缩减和普通丢包；在IncreaseWindow路径中继续检查ECN状态。RTT膨胀用于参数调节和慢启动退出，不单独触发乘性降窗。

（2）\textbf{启发式反馈自适应的融合控制}：BDP由滑动时间窗口内的累计确认字节差分估计交付速率，再经容量40的窗口最大值滤波并乘以最小RTT得到。控制器根据最小RTT感知的RTT膨胀阈值、反馈值相对慢速基线的偏移和连续增长趋势调节每流$\alpha\in[0.85,1.30]$，在慢启动与拥塞避免阶段以有界步长调整窗口。

（3）\textbf{稳定性与安全约束}：不同拥塞类型分别采用丢包0.70、ECN 0.75和超时0.50的窗口保留因子；连续缩减保护、降窗后4个ACK事件冻结、窗口上下界以及CA\_LOSS状态下的陈旧动作作废共同限制窗口过度收缩、快速反弹和陈旧决策覆盖。上述控制回路经ns3-gym/OpenGym与ZeroMQ同步执行，当前实验采用单一随机种子，其跨种子稳定性仍需重复实验检验。
